\scanpage{12}
(iv) Supposons \deleted{d'abord} $S$ \deleted{régulier}\inserted{normal}, \inserted{donc aussi tous les schémas écrits, \deleted{\unsure{sauf $G'$}}} \deleted{Alors $\mathrm{Pic}(\mathscr B)\xrightarrow{\sim}\mathrm{Pic}(\mathscr B^{\mathrm{reg}})$} $\mathrm{Pic}(\mathscr B)\longrightarrow\mathrm{Pic}(\mathscr B^{\mathrm{reg}})$ injectif résulte \inserted{\unsure{notions|questions|considérations} de profondeur et} du fait que $\mathscr B-\mathscr B^{\mathrm{reg}}$ est de codim $\geq 2$, \deleted{\unsure{et $\mathrm{Pic}(G'$}} \inserted{et $\mathrm{Pic}(D)\longrightarrow\mathrm{Pic}(\mathscr B^{\mathrm{reg}})$} résulte de considération fibre par fibre.

Le fait que $\mathrm{Pic}(D)\longrightarrow\mathrm{Pic}(\mathscr B)$ \deleted{soit un isom.}\inserted{\uline{surjectifs}} résulte de EGA $\underline{\mathrm{IV}}$ 21.4.13 (\unsure{Err|EGA} $\underline{\mathrm{IV}}$ 53), compte tenu du

\uline{Lemme}\quad $G$ groupe \deleted{\unsure{réductif (resp.)}}\inserted{semi-simple} simplement connexe sur corps $k$, $B$ un Borel, alors $\mathrm{Pic}(B)\deleted{\unsure{=\mathrm{Pic}(T)}}\inserted{=\mathrm{Pic}(B^{\mathrm{reg}})}=\mathrm{Pic}(T)=H^1(k,\underline M)$ \deleted{\unsure{si $k$}} \inserted{($T=B/B_u$, $\underline M=\widehat T$)}

Par loc. cit, utilisant \deleted{la relation} $B\longrightarrow T=B/B_u$, on est ramené au cas où $\mathrm{Pic}(T)=0$ (NB $B_u$ se dévisse en groupes additifs).

\margin{On a $=$ $\mathrm{Pic}(B)\simeq\mathrm{Pic}(T)$}

Or on sait que $T=\prod_{k'/k} T'$, $T'$ un tore sur une \deleted{\unsure{algèbre}}\inserted{extension} étale \inserted{galoisienne (gpe $\Gamma$)}, donc on a \deleted{\unsure{une équivalence de cat.|équiv. de cat.}}
\[
\deleted{\underline{\mathrm{Inv}}(T)\simeq\underline{\mathrm{Inv}}(T')\quad\text{donc}}
\]
\[
\deleted{\mathrm{Pic}(T)}\quad
0\longrightarrow H^1\bigl(\Gamma,\underbrace{H^0(T',\underline O_{T'}^{\times})}_{M\times k'^{*}}\bigr)
\longrightarrow\mathrm{Pic}(T)\longrightarrow H^0\bigl(\Gamma,\underbrace{H^1(T',\underline O_{T'}^{\times})}_{=0}\bigr)
\]
\deleted{\unsure{partie adjointe est nulle} \unsure{au sens entier}}

donc on trouve
\[
\mathrm{Pic}(T)=H^1(\Gamma,M)\times\underset{\substack{=0\\\mathrm{H.}\,90}}{H^1(\Gamma,k'^{\times})}=H^1(\Gamma,M)=H^1(k,\underline M)
\]
Si $P$ est tordu, cela peut être $\ne 0$ ; si $B$ est la fibre générique de $\mathscr B$ sur $D_{\unsure{s|S|{}}}$ ($S$ irréd), alors cela implique donc que l'hom $\mathrm{Pic}(D)\longrightarrow\mathrm{Pic}(\mathscr B)$ n'est pas surjectif. L'injectivité est claire, car $\mathscr B$ a une section sur $D$ (la section nulle). En faisceautisant pour la topologie étale, on trouve donc le diagramme d'isoms sur les $\underline{\mathrm{Pic}}$. / Comme $\mathscr B\longrightarrow D$ est affine, et a une section, il s'ensuit qu'un faisceau inv. sur $D$ est ample ssi son image inverse sur $\mathscr B$ l'est.
\margin{indép. de l'hypothèse sur $S$ !} % slanted

\deleted{Que peut-on faire dans le cas non régulier ? $\Rightarrow\underline{\mathrm{Pic}}$ \unsure{$\subset$} a) L'argument établissant $\mathrm{Pic}(D)\xrightarrow{\sim}\mathrm{Pic}(\mathscr B)$ \unsure{marche|et} dès que $S$ est \uline{normal}. De $=$ pour $\mathrm{Pic}(D)\longrightarrow\mathrm{Pic}(\mathscr B^{\mathrm{reg}})$. L'injectivité t\textsuperscript{js}. \unsure{directement|décentrement|décentement|évidemment} !}
% archive sheet number: 11
